
\section{Introducción}

En el siguiente informe se detalla 

\section{Objetivo}

\section{Marco teórico}

% Qué son los dispositivos de medición, decisión y acción?

Operaciones básicas en todo sistema de control:

\begin{itemize}
    \item Medición (M): de la variable que se controla. Se hace generalmente mediante la combinación sensor-transmisor.
    \item Decisión (D): en base a la medición el controlador decide que hacer para mantener a la variable en el valor deseado.
    \item Acción (A): como resultado de la decisión se debe efectuar una acción en el sstema, generalmente realizada por el elemento final de control.
\end{itemize}

\section{Desarrollo}

\subsection{Consigna 1}

Para cada uno de los siguientes controles automáticos encontrados en la vida
diaria, identificar los dispositivos que realizan funciones de medición (M), decisión
(D) y acción (A). Dibujar un diagrama de instrumentos y cañerías (P\&ID) utilizando
la norma ISA S5.1 y determinar si el sistema de control es realimentado o por
acción precalculada, a lazo abierto o a lazo cerrado.

\begin{enumerate}
    \item Aire acondicionado doméstico.

Dispositivos:

\begin{itemize}
    \item Medición: Termistor que compone el termostato.
    \item Decisión: Termostato. Es el encargado deencender y apagar la unidad y mantener la temperatura prestablecida
    \item Acción: Condensador
\end{itemize}



    \item Horno.
    \item Tostadora.
    \item Sistema de sprinklers contra incendio.
    \item Control automático de velocidad de crucero en automóviles.
    \item Heladeras domésticas.
    \item Puertas automáticas en los accesos de centros comerciales.
    \item Sistemas de control automático de temperatura para bañeras.
\end{enumerate}

\section{Fórmulas TP2}

\subsection{Workshop 4}

$$ K_{c}= -2.04 $$

$$ \tau_{I}= 0.727 $$

\subsubsection{point 5}

$$ K_{c}= 2.04 \frac{\%}{°C} \, \, \, \, \, \tau_{I}= 0.75 \, min \, \, \, \, \, \tau_{D}= 0.295 \, min$$

\subsubsection{point 7}

$$ K_{c}= -2.04 $$


\subsection{Workshop 5}

$$ y_{setpoint}= 90.2\% $$

\subsubsection{point 5}

$$ K_{c}= 0.935  \, \, \, \,  \tau_{I}= 43.86 $$

\subsubsection{point 8}

\subsubsection{case 1}

\subsection{point 6}

$$ K_{c}= -0.4858  \, \, \, \, \tau_{I}= 1.28 \, \, \, \, \theta_{D}= 0.579$$


% Referencias
% https://www.climamania.com/blog/partes-aire-acondicionado/?srsltid=AfmBOorn9Gjoa8HXR36uT12_65nbObxzSEUdYVKTOccCmrhNlabAr3li